\mxunit{seconde-s05}{Bilan de trimestre et remédiation}\mxsubtitle{16--20 novembre 2026}
\section{Bilan individuel}
\begin{method}[Consigne]
Chercher sans le cours pendant 25 minutes, puis vérifier avec les méthodes du livret. Pour chaque question, noter : réussi seul, réussi avec aide, ou à reprendre. Justifier les réponses.
\end{method}
\begin{example}[Bilan des séquences 1 à 4]\begin{enumerate}
\item Rendre $90/126$ irréductible. Démontrer que la somme de deux entiers impairs est paire.
\item Avec $A\pt{-1;2}$ et $B\pt{3;-2}$, calculer $\vect{AB}$, $AB$ et le milieu de $[AB]$ dans un repère orthonormé.
\item Déterminer une équation de $(AB)$ ; $C\pt{5;-4}$ appartient-il à cette droite ?
\item Calculer $\intff{-2}{4}\cap\intof{1}{7}$ et $\intff{-2}{4}\cup\intof{1}{7}$.
\item Résoudre $|x-1|\leq4$ et expliquer pourquoi $1/3\notin\D$.
\end{enumerate}\end{example}
\mxanswerspace[9]{1. $\pgcd(90,126)=18$, donc $90/126=5/7$. Si $u=2k+1$ et $v=2\ell+1$, alors $u+v=2(k+\ell+1)$ est pair.\\2. $\vect{AB}=\pt{4;-4}$ ; $AB=4\sqrt2$ ; milieu $I\pt{1;0}$.\\3. $y=-x+1$. Pour $C$, $-4=-5+1$ : il appartient à $(AB)$.\\4. Intersection $\intof{1}{4}$ ; réunion $\intff{-2}{7}$.\\5. $x\in\intff{-3}{5}$. Une écriture $1/3=a/10^n$ imposerait que $10^n$ soit divisible par $3$, ce qui est impossible.}
\begin{teachernote}Barème possible sur 20 : 4 points par question. Repérer séparément maîtrise du calcul, choix de méthode et qualité de la justification.\end{teachernote}
\mxsection{Corriger une erreur et reprendre une méthode}
\begin{example}[Analyse de productions]Corriger : (a) « $\vect{AB}=\pt{x_A-x_B;y_A-y_B}$ » ; (b) « Le vecteur nul dirige toutes les droites » ; (c) « $|x-2|\leq3$ donne $x\in\intff{2}{3}$ » ; (d) « $1/3\approx\num{0,333}$, donc $1/3$ est décimal ».\end{example}
\mxanswerspace[5]{(a) Soustraire origine à extrémité : $\pt{x_B-x_A;y_B-y_A}$. (b) Un directeur doit être non nul. (c) La distance à $2$ est au plus $3$ : $x\in\intff{-1}{5}$. (d) $\num{0,333}$ est différent de $1/3$ : $3\times\num{0,333}=\num{0,999}\ne1$. Une approximation décimale ne détermine pas la nature du nombre approché.}
\begin{example}[Parcours de reprise]\begin{enumerate}
\item Réduire $72/108$, puis $\frac34-\frac16$.
\item Avec $P\pt{1;1}$ et $Q\pt{5;3}$, déterminer le milieu, un directeur et une équation de $(PQ)$.
\item Déterminer $\intof{-\infty}{3}\cap\intfo{-1}{+\infty}$, puis résoudre $|x+2|<1$.
\end{enumerate}\end{example}
\mxanswerspace[5]{1. $2/3$ ; $7/12$. 2. Milieu $\pt{3;2}$, directeur $\pt{2;1}$, équation $x-2y+1=0$. 3. Intersection $\intff{-1}{3}$ ; solution de l'inéquation $\intoo{-3}{-1}$.}
\mxsection{Réinvestir et s'autoévaluer}
\begin{example}[Problème de synthèse]Deux déplacements sont représentés par $\vec u\pt{6;4}$ et $\vec v\pt{9;6}$.
\begin{enumerate}\item Montrer qu'ils sont colinéaires et trouver le coefficient reliant $\vec v$ à $\vec u$.
\item Le départ est $A\pt{-2;1}$. Déterminer les arrivées $B$ et $C$ associées aux déplacements $\vec u$ et $\vec v$.
\item Trouver une équation de la droite passant par ces trois points.
\item Un point $M\pt{x;y}$ de cette droite doit vérifier $|x-4|\leq3$. Donner l'intervalle des abscisses possibles.
\end{enumerate}\end{example}
\mxanswerspace[7]{1. $\vec v=\frac32\vec u$. 2. $B\pt{4;5}$ et $C\pt{7;7}$. 3. Un directeur est $\pt{3;2}$ ; l'équation est $2x-3y+7=0$, vérifiée par les trois points. 4. $|x-4|\leq3\iff x\in\intff{1}{7}$.}
\begin{keypoint}[Mon bilan]
Ma priorité : \dotfill\par L'erreur à éviter : \dotfill\par La méthode à réutiliser : \dotfill\par Un exercice à refaire sans aide : \dotfill
\end{keypoint}
\begin{teachernote}Valider la reprise par un exercice analogue réalisé sans aide après un délai. À l'oral, demander une phrase reliant calcul et conclusion géométrique.\end{teachernote}
